%============================================================
%  LEX v3.0 — Template Surat Pengantar Submission Jurnal
%  Gunakan file ini saat mengirimkan naskah ke jurnal
%  Isi setiap placeholder [ ] sebelum dikirim
%============================================================

\documentclass[12pt,a4paper]{article}

%--- Paket utama ---
\usepackage[bahasa]{babel}
\usepackage[utf8]{inputenc}
\usepackage[T1]{fontenc}
\usepackage{times}
\usepackage[top=3cm, bottom=3cm, left=3cm, right=2.5cm]{geometry}
\usepackage{setspace}
\singlespacing
\usepackage{parskip}
\usepackage{hyperref}
\hypersetup{colorlinks=true, linkcolor=black, urlcolor=blue}

%============================================================
\begin{document}

%------------------------------------------------------------
%  KANAN ATAS: Kota dan Tanggal
%------------------------------------------------------------
\begin{flushright}
  [Kota], \today
\end{flushright}

\bigskip

%------------------------------------------------------------
%  TUJUAN SURAT
%------------------------------------------------------------
Kepada Yth.\\
Editor-in-Chief / Dewan Redaksi\\
\textbf{[Nama Jurnal]}\\
[Institusi Penerbit]\\
[Alamat / URL Jurnal]

\bigskip

%------------------------------------------------------------
%  SALAM PEMBUKA
%------------------------------------------------------------
Dengan hormat,

\bigskip

%------------------------------------------------------------
%  PARAGRAF 1 — Identifikasi Naskah
%------------------------------------------------------------
Kami mengajukan naskah dengan judul \textbf{``[Judul Artikel]''} untuk dipertimbangkan
diterbitkan pada \textit{[Nama Jurnal]} (\textit{[ISSN/e-ISSN]}). Naskah ini
% Hapus salah satu pilihan di bawah (artikel penelitian ATAU artikel konseptual)
merupakan artikel penelitian/artikel konseptual
di bidang hukum [bidang spesifik, misal: hukum tata negara / hukum pidana / hukum internasional],
dengan panjang sekitar [X] kata (tidak termasuk abstrak dan daftar pustaka).

\bigskip

%------------------------------------------------------------
%  PARAGRAF 2 — Ringkasan Kontribusi (Novelty)
%  Copilot: gunakan /research-gap untuk mengisi bagian ini
%------------------------------------------------------------
Artikel ini menganalisis [masalah hukum utama] dengan menggunakan kerangka Asas-Teori-Doktrin
(ATD). Kontribusi utama naskah ini adalah: pertama, [novelty 1]; kedua, [novelty 2]; dan
ketiga, [novelty 3]. Temuan ini belum dibahas secara memadai dalam literatur yang ada,
sebagaimana ditunjukkan dalam kajian pustaka kami.

\bigskip

%------------------------------------------------------------
%  PARAGRAF 3 — Pernyataan Etika dan Eksklusivitas
%------------------------------------------------------------
Kami menyatakan bahwa:
\begin{enumerate}
  \item Naskah ini merupakan karya orisinal, bebas plagiarisme, dan belum pernah
        diterbitkan atau sedang dalam proses review di jurnal lain.
  \item Seluruh penulis telah menyetujui isi naskah dan urutan nama penulis.
  \item Tidak terdapat konflik kepentingan yang perlu diungkapkan.
  \item Penelitian ini [tidak memerlukan / telah memperoleh] persetujuan etik
        dari [nama komite etik, jika ada].
\end{enumerate}

\bigskip

%------------------------------------------------------------
%  PARAGRAF 4 — Saran Reviewer (opsional, tapi dianjurkan)
%------------------------------------------------------------
Kami menyarankan beberapa nama yang kompeten untuk menelaah naskah ini (opsional):
\begin{itemize}
  \item [Nama Reviewer 1] — [Institusi] (\texttt{[email@institusi.ac.id]})
  \item [Nama Reviewer 2] — [Institusi] (\texttt{[email@institusi.ac.id]})
\end{itemize}

\bigskip

%------------------------------------------------------------
%  PARAGRAF 5 — Penutup
%------------------------------------------------------------
Kami berharap naskah ini memenuhi standar keilmuan \textit{[Nama Jurnal]} dan layak untuk
ditelaah lebih lanjut. Kami siap memberikan klarifikasi ataupun revisi sesuai masukan
dewan redaksi dan reviewer.

\bigskip

Hormat kami,

\vspace{2em}

\noindent
[Nama Penulis Korespondensi]\\
[Jabatan / Program Studi]\\
[Nama Institusi]\\
Telp./WA: [Nomor Telepon]\\
Surel: \href{mailto:[email@institusi.ac.id]}{[email@institusi.ac.id]}\\
ORCID: \href{https://orcid.org/[ORCID-ID]}{https://orcid.org/[ORCID-ID]}

\bigskip
\hrule
\medskip
{\footnotesize
\textbf{Penulis Lain:}\\
[Nama Penulis 2] — [Institusi] (\texttt{[email]})\\
[Nama Penulis 3] — [Institusi] (\texttt{[email]})
}

\end{document}
